% =====================================================================
%  ANNEXES
%  Qualité rédactionnelle équivalente au texte principal.
%  Chaque annexe doit être référencée dans le corps du rapport.
% =====================================================================
\appendix
\chapter*{Annexes}
\addcontentsline{toc}{chapter}{Annexes}
\markboth{Annexes}{Annexes}

% Numérotation des annexes : A, B, C...
\renewcommand{\thesection}{\Alph{section}}
\setcounter{section}{0}
\titleformat{\section}
  {\normalfont\normalsize\bfseries\color{vinciblue}}{Annexe~\thesection}{1em}{}

\section{Évolution de l'interface de \csbag}
\label{ann:v1v2}

Les captures ci-dessous illustrent l'amélioration de l'ergonomie décrite en
section~\ref{sec:existant}. Pour chaque page, la version~1 (à gauche) est mise
en regard de la version~2 (à droite).

\subsection*{Page d'accueil}

\begin{figure}[H]
  \centering
  \begin{subfigure}[t]{0.48\textwidth}
    \centering
    \screenshot[\linewidth]{csbag_v1_accueil.png}
    \caption{Version~1~: liste sans recherche ni filtre.}
    \label{fig:v1-accueil}
  \end{subfigure}
  \hfill
  \begin{subfigure}[t]{0.48\textwidth}
    \centering
    \screenshot[\linewidth]{csbag_v2_accueil.png}
    \caption{Version~2~: barre de recherche et filtre d'équipe.}
    \label{fig:v2-accueil}
  \end{subfigure}
  \caption{Page d'accueil~: avant et après refonte.}
  \label{fig:ann-accueil}
\end{figure}

\subsection*{Page « Tout voir »}

\begin{figure}[H]
  \centering
  \begin{subfigure}[t]{0.48\textwidth}
    \centering
    \screenshot[\linewidth]{csbag_v1_toutvoir.png}
    \caption{Version~1~: vue sans filtres ni détail dépliable.}
    \label{fig:v1-toutvoir}
  \end{subfigure}
  \hfill
  \begin{subfigure}[t]{0.48\textwidth}
    \centering
    \screenshot[\linewidth]{csbag_v2_toutvoir.png}
    \caption{Version~2~: filtres par équipe, référent et framework.}
    \label{fig:v2-toutvoir}
  \end{subfigure}
  \caption{Page « Tout voir »~: avant et après refonte.}
  \label{fig:ann-toutvoir}
\end{figure}

\subsection*{Détail d'une application}

\begin{figure}[H]
  \centering
  \begin{subfigure}[t]{0.48\textwidth}
    \centering
    \screenshot[\linewidth]{csbag_v1_detail_app.png}
    \caption{Version~1~: affichage des modules et ressources.}
    \label{fig:v1-detail-app}
  \end{subfigure}
  \hfill
  \begin{subfigure}[t]{0.48\textwidth}
    \centering
    \screenshot[\linewidth]{csbag_v2_detail_app.png}
    \caption{Version~2~: modules enrichis, référents avec fonction.}
    \label{fig:v2-detail-app}
  \end{subfigure}
  \caption{Détail d'une application~: avant et après refonte.}
  \label{fig:ann-detail-app}
\end{figure}

\subsection*{Détail d'un module}

\begin{figure}[H]
  \centering
  \begin{subfigure}[t]{0.48\textwidth}
    \centering
    \screenshot[\linewidth]{csbag_v1_detail_module.png}
    \caption{Version~1~: ressources en texte libre.}
    \label{fig:v1-detail-module}
  \end{subfigure}
  \hfill
  \begin{subfigure}[t]{0.48\textwidth}
    \centering
    \screenshot[\linewidth]{csbag_v2_detail_module.png}
    \caption{Version~2~: ressources normalisées et triables.}
    \label{fig:v2-detail-module}
  \end{subfigure}
  \caption{Détail d'un module~: avant et après refonte.}
  \label{fig:ann-detail-module}
\end{figure}

\subsection*{Modification d'une application}

\begin{figure}[H]
  \centering
  \begin{subfigure}[t]{0.48\textwidth}
    \centering
    \screenshot[\linewidth]{csbag_v1_modif_app.png}
    \caption{Version~1~: formulaire d'édition.}
    \label{fig:v1-modif-app}
  \end{subfigure}
  \hfill
  \begin{subfigure}[t]{0.48\textwidth}
    \centering
    \screenshot[\linewidth]{csbag_v2_modif_app.png}
    \caption{Version~2~: formulaire enrichi avec auto-complétion.}
    \label{fig:v2-modif-app}
  \end{subfigure}
  \caption{Modification d'une application~: avant et après refonte.}
  \label{fig:ann-modif-app}
\end{figure}

\subsection*{Modification d'un module}

\begin{figure}[H]
  \centering
  \begin{subfigure}[t]{0.48\textwidth}
    \centering
    \screenshot[\linewidth]{csbag_v1_modif_module.png}
    \caption{Version~1~: saisie libre des ressources.}
    \label{fig:v1-modif-module}
  \end{subfigure}
  \hfill
  \begin{subfigure}[t]{0.48\textwidth}
    \centering
    \screenshot[\linewidth]{csbag_v2_modif_module.png}
    \caption{Version~2~: sélection depuis les référentiels normalisés.}
    \label{fig:v2-modif-module}
  \end{subfigure}
  \caption{Modification d'un module~: avant et après refonte.}
  \label{fig:ann-modif-module}
\end{figure}

\subsection*{Certificats}

\begin{figure}[H]
  \centering
  \begin{subfigure}[t]{0.48\textwidth}
    \centering
    \screenshot[\linewidth]{csbag_v1_certificats.png}
    \caption{Version~1~: liste basique des certificats.}
    \label{fig:v1-certificats}
  \end{subfigure}
  \hfill
  \begin{subfigure}[t]{0.48\textwidth}
    \centering
    \screenshot[\linewidth]{csbag_v2_certificats.png}
    \caption{Version~2~: empreinte, recherche et alertes d'expiration.}
    \label{fig:v2-certificats}
  \end{subfigure}
  \caption{Gestion des certificats~: avant et après refonte.}
  \label{fig:ann-certificats}
\end{figure}

\section{Principaux écrans ajoutés dans \csbag v2}
\label{ann:ecrans}

Cette annexe présente les écrans qui n'existaient pas dans la version~1, décrits
en section~\ref{sec:fonctionnel}.

\begin{figure}[H]
  \centering
  \begin{subfigure}[t]{0.48\textwidth}
    \centering
    \screenshot[\linewidth]{csbag_historique_actions.png}
    \caption{Historique des actions~: traçabilité avec affichage hiérarchisé.}
    \label{fig:ann-historique}
  \end{subfigure}
  \hfill
  \begin{subfigure}[t]{0.48\textwidth}
    \centering
    \screenshot[\linewidth]{csbag_admin_referentiels.png}
    \caption{Administration d'un référentiel issu de la normalisation.}
    \label{fig:ann-admin}
  \end{subfigure}
  \caption{Écrans ajoutés dans \csbag v2.}
  \label{fig:ann-ecrans}
\end{figure}

\section{Suivi des travaux dans Wrike}
\label{ann:wrike}

Cette annexe illustre l'outil de gestion de projet évoqué en
section~\ref{sec:organisation}. Wrike est organisé en espaces qui reprennent la
structure de l'organisation~: l'espace CS Group regroupe les squads M365 et
MyVIEW, ainsi que les tableaux de suivi transverses du service, parmi lesquels
le suivi du support et la feuille de route. Les projets internes du squad, dont
\csbag, y disposent chacun de leur propre dossier.

Chaque lot fonctionnel fait l'objet d'un projet distinct, suivi sur un tableau
agile dont les colonnes correspondent aux états d'avancement, de la tâche en
cours à la tâche terminée, en attente ou annulée. Les tâches sont formulées à
la maille de la fonctionnalité et décomposées en sous-tâches, dont le nombre
apparaît sur chaque carte. Ce découpage rend visible, à tout moment et pour
l'ensemble de l'équipe, ce qui reste à produire et ce qui a été livré, et sert
de support aux points d'avancement hebdomadaires.

\begin{figure}[H]
  \centering
  \screenshot{wrike_lot1.png}
  \caption{Tableau agile du premier lot de \csbag v2 dans Wrike.}
  \label{fig:ann-wrike}
\end{figure}
